% UNINTEGRATED ARITHMETIC15 FOUNDATION PART, 25 October2026.
% Same existing arithmetic paper family; source preparation only.
% First batch: accepted general dihedral3/19 theorem and D14 consumer.
% DRAFT pending independent and coordinating review; no execution/compilation.
\section{A strengthened ordinary dihedral comparison}
\label{app:arithdihedral}

This section supplies the ordinary rank-two comparison needed beside
\(D_{14}\).  Its component theorem is general real geometry.
Arithmeticity enters only when the already established arithmetic
\(D_{14}\) prefix bound is used in the final corollary.
Throughout, \(p\) is the ORIGINAL vertex farthest from one strict
interior point.  The theorem is not asserted at an arbitrary
neighbour with a weakened comparison factor.

\begin{proposition}\label{prop:ordinarydihedralthreenineteenths}
Let a compact hyperbolic Coxeter polytope have an actual farthest
vertex \(p\), and write \(z=E(p-x)\), \(\|x\|<1\).
Let one WHOLE component of the incident finite diagram be
\(I_2(m)\), \(m\ge4\), and let \(e\) be the squared norm of the
projection of this SAME \(x\) onto that component.
If \(e\le3/19\), then \(m=7\) and one actual opposite wall
closes both component directions, has a single old contact of
label three there, and has no contact with any other old component.
This is precisely the actual \(O_7\) profile.
Consequently every ordinary \(I_2(m)\), \(m\ge4\), has energy
strictly greater than \(3/19\), with arbitrary finite spectators.
\end{proposition}

\begin{proof}
Use outward unit incident roots, finite Gram \(B\), inverse
\(C=B^{-1}\ge0\), duals \(\omega_i\), and unit inward dual
directions \(-u_i=-\omega_i/\sqrt{C_{ii}}\).
For an actual outgoing edge of length \(L_i\), maximality at
the original \(p\) gives
\[
 x\cdot u_i\ge\tau_i=\tanh(L_i/2)>0,\qquad
 r_i=\tanh L_i={2\tau_i\over1+\tau_i^2}.
\]
The old dihedral inverse is
\[
 C_I=d\begin{pmatrix}1&k\\k&1\end{pmatrix},\qquad
 k=\cos(\pi/m)\ge1/\sqrt2,\quad d=(1-k^2)^{-1}.
\]
Thus each \(\tau_i^2\le e\), and the increasing function
\(4t/(1+t)^2\) on \(0\le t<1\) gives
\[
                         r_i^2\le57/121=:r .
\tag{Dih1}\label{eq:arithdihedralr}
\]

We recall the actual-wall calculation forcing a common opposite.
For an outside actual wall \(c=Ap-v\), let
\(\beta_j=-c\cdot n_j\ge0\) and put
\[
 T_i(c)={(C\beta)_i\over A\sqrt{C_{ii}}}.
\]
For the actual opposite \(c_i\) of edge \(i\),
\(T_i(c_i)=1/r_i\).
The matrix
\(C-C_{\cdot j}C_{j\cdot}/C_{jj}\)
is entrywise nonnegative: its nonzero block is the inverse of
the positive principal Gram obtained by deleting \(j\).
Actual acuteness between two distinct opposites therefore gives
\(T_j(c_i)\le r_j\), while that same nonnegative matrix gives
\(T_j(c_i)\ge k/r_i\) within this old component.
Hence distinct opposites imply \(r_ir_j\ge k\).
But~\eqref{eq:arithdihedralr} gives
\(r_ir_j\le57/121<1/\sqrt2\le k\).
One actual wall \(c\) therefore closes BOTH directions.

Write its full decomposition as
\[
 c=Ap-v_I-v_S,\quad h=A^2>0,\quad
 b=-c\cdot n_1\ge0,\quad g=-c\cdot n_2\ge0,\quad
 Q=\|v_S\|^2\ge0 .
\]
Here \(S\) is the entire orthogonal spectator forest.
Every positive contact is a native magnitude at least \(1/2\);
none of its finite, parallel or dotted possibilities is omitted
before the actual endpoint constraints are imposed.
At least one of \(b,g\) is positive.

Suppose both are positive.
At the actual endpoint releasing direction one, the retained
white Gram has residual \(\delta=1-Q-g^2>0\).
Its actual edge formula is
\[
 \cosh L_1={\sqrt d(b+kg)\over\sqrt{1-Q-g^2}}
 \ge{\sqrt d(1+k)\over\sqrt3}
 ={\,\cot(\pi/(2m))\over\sqrt3}
 \ge{1+\sqrt2\over\sqrt3}>{11\over8}.
\]
The last comparison is \(128\sqrt2>171\), whose positive
squared margin is \(32768-29241\).
It follows that
\[
 e\ge\tanh^2(L_1/2)>
 {11/8-1\over11/8+1}={3\over19},
\]
a contradiction.  Hence exactly one old contact \(\beta>0\)
remains.

We need one whole-spectator fact at small time.
If \(h<1/4\), then \(c\) cannot contact any foreign component.
Indeed a positive foreign contact gives a normalized derivative
\[
 {(C_S\beta)_j\over\sqrt{(C_S)_{jj}}}\ge1/2 .
\]
This uses the full spectator inverse, with nonnegative entries
and diagonal at least one.
Its ratio to \(\sqrt h\) exceeds one.
The actual-wall shadow inequality therefore forces this SAME
\(c\) to be the opposite of that foreign edge as well.

For clarity, the required orthogonal obstruction is specific to
the original farthest vertex.  Suppose a wall closed orthogonal
dual directions \(u_i,u_j\).
Clear it positively off all other incident roots \(W\):
\[
 c_F=c+\sum_{s\in W}\alpha_sn_s,\qquad
 \alpha=B_W^{-1}(-c\cdot n_W)\ge0.
\]
Strict interiority gives \(z\cdot c_F<0\).
In the retained Lorentz three-space its exact form is
\[
 c_F=A p-(A/r_i)u_i-(A/r_j)u_j .
\]
Therefore
\[
 {x\cdot u_i\over r_i}+{x\cdot u_j\over r_j}<1.
\]
Original maximality gives the opposite bound
\(\tau_i/r_i+\tau_j/r_j
 =1+(\tau_i^2+\tau_j^2)/2>1\).
This proves the obstruction, including arbitrary contacts in \(W\).
Thus \(h<1/4\) really forces \(Q=0\).

Put \(q=d\beta^2\).
The smaller normalized derivative square is \(k^2q\).
The common-wall edge formula and~\eqref{eq:arithdihedralr} imply
\[
 q-1\le h=q+Q-1\le rk^2q .
\]
Substituting \(k^2=1-1/d\) and \(r=57/121\) gives the exact
strengthened restriction
\[
                         \beta^2(64d+57)\le121 .
\tag{Dih2}\label{eq:arithdihedralcut}
\]
This is derived at the \(3/19\) threshold; it is not the older
\(1/7\) inequality with an unchanged coefficient.

Since \(\beta^2\ge1/4\), this forces \(d\le427/64\).
For every old \(m\ge8\),
\(d\ge4+2\sqrt2>34/5>427/64\), impossible.
The complete remaining old alphabet is \(4,5,6,7\).
At old \(m=5\),
\(d=2+2/\sqrt5>185/64\):
equivalently \(128>57\sqrt5\), with squared margin139.
Since \(d\) increases with \(m\), if
\(\beta\ge1/\sqrt2\) at any of old \(5,6,7\),
\eqref{eq:arithdihedralcut} would give the opposite
\(d\le185/64\).
Thus all three rows have exactly \(\beta=1/2\).
Every larger finite contact, parallel and dotted contact has
been excluded by the native jump, rather than a sampled label list.

At old four, \(d=2\) and
\(\beta^2\le121/185<\cos^2(\pi/5)\).
The strict comparison is \(413<185\sqrt5\), with squared
margin556.  Only contacts three and four remain, giving
\(q=1/2\) or \(1\).
Consequently \(0<h\le57/242<1/4\).
The small-time argument just proved gives \(Q=0\), whereas
\(h=q-1\le0\), a contradiction.

At old five only \(\beta=1/2\) remains.
Since \(d<3\), \(q<3/4\) and
\(0<h\le r(d-1)/4<57/242<1/4\).
Again \(Q=0\), and \(h=q-1<0\).

At old six, \(\beta=1/2\), \(d=4\), \(q=1\).
If \(Q=0\), then \(h=0\), impossible.
For any \(Q>0\), the actual positive sum
\[
                          \nu=c+2n_1+\sqrt3n_2
\]
is a nonzero future null: all three selected products vanish.
The coefficient order is for OLD6/NEW3, since
\(c\cdot n_1=-1/2\), \(n_1\cdot n_2=-\sqrt3/2\),
\(c\cdot n_2=0\).
Every other actual wall pairs nonpositively termwise with \(\nu\).
It is a feasible ideal ray, contrary to compactness.
No foreign contact was discarded in this affine exclusion.

At old seven only \(\beta=1/2\) remains, and
\[
 0\le Q=1+h-d/4
 \le{427-64d\over484}<{107\over484}<{1\over4},
\]
using \(d_7>5\).
Every positive spectator cost is at least \(1/4\), by inverse
positivity and the native contact lower bound.  Thus \(Q=0\).
We have exactly the ACTUAL isolated-new-three wall in old
\(I_2(7)\), closing both directions: the stated \(O_7\).
Its time square is \(d_7/4-1\in(0,1/3)\);
the bounds \(5<d_7<16/3\) follow from
\(6/5<2\cos(2\pi/7)<5/4\).
Its separate inherited energy floor \(>9/100\) is unchanged.

All weak inequalities included \(e=3/19\), so no ordinary equality
case is omitted.  The proof has retained every finite spectator
component, all its contacts and the same original \(x\).
No arithmetic hypothesis or weakened-neighbour comparison was used.
\end{proof}

\begin{corollary}\label{cor:D14ordinarydihedral}
At the original farthest vertex in arithmetic dimension sixteen,
a whole \(D_{14}\) cannot have an ordinary irreducible
rank-two spectator.
\end{corollary}
\begin{proof}
The fresh true-rank-\(14\) argument in the proof of
Proposition~\ref{prop:f16} gives
\[
                         e_D>.84816={5301\over6250}
\]
for its actual whole rank-two spectator.
Its bootstrap uses only the independently available spectator
floor \(>.09\), so it does not depend on the present strengthening.
For an ordinary \(I_2(m)\), \(m\ge4\), the proposition gives
\[
 \|x\|^2=e_D+e_S>
 {5301\over6250}+{3\over19}
 =1+{719\over118750}>1 .
\]
For ordinary \(A_2\), use its geometric companion bound
\(319/2000\), which exceeds \(3/19\) by \(61/38000\).
Both contradict \(\|x\|<1\).
The actual typed \(N_7/O_7\) alternatives retain their separate
costs and arguments; they have not been charged as ordinary.
\end{proof}

% Next batches may append reviewed arithmetic proof bodies here.
% Preserve this exact first-batch source before any such expansion.

\section{Arithmetic chain windows and simultaneous actual feet}
\label{app:arithmeticchain}

The endpoint arguments below share one arithmetic chain calculation.
We prove it before specializing to the whole \(A_{16}\) and
\(A_{17}\) cases.  The geometric finite-white classification is the
one in the general companion; the arithmetic restrictions and the
simultaneous treatment of isolated feet are given here.
No component dimension is obtained by deleting a node from a
previous ambient theorem.

Let \(N=16,17,18\), let the original incident forest contain the
WHOLE component \(A_{N-1}\), and retain all finite spectators.
Use its zero-sum \(\mathbb R^N\) frame and unit old roots
\(n_i=(e_i-e_{i+1})/\sqrt2\).  Then
\[
 D_i=\omega_i^2={2i(N-i)\over N},\qquad
 C_{ij}=\omega_i\cdot\omega_j={2i(N-j)\over N}\quad(i\le j).
\]
All heights \(h_i=x_A\cdot\omega_i\) concern the SAME original
farthest vector.  Set
\[
             \ell_{16}=.601,\qquad\ell_{17}=\ell_{18}=.6.
\]

\subsection{The complete central inventory at the new losses}
\label{app:arithchaininventory}

For \(5\le i\le N-5\), call an index weak when
\(h_i\le\sqrt{D_i}-\ell_N\).
We claim that every central constraint is either strong, an
isolated-within-\(A\) cap with its actual cleared parameter
\(0<u<1\), or a pure golden end-helper.
The last alternative can occur at at most ONE central index.

First establish the full inventory, without assuming a contrary
component energy.  Deleting \(i\) retains two old paths of ranks
at least four.  Their complete one-node finite extensions are
isolation; one admissible end, \(D\)- or \(E\)-position three;
one terminal four; or two endpoint threes, one in each path.
Two contacts in a single path make a cycle.
A connected two-path completion has rank \(N-1>8\), so finite
classification forces the ordinary path bridge, with no extra
finite foreign arm.

The zero-own repair depends only on side order \(m\ge5\).
At a far-end three the time square is less than its full positive
foreign cost; the next-to-far three has foreign remainder \(<2/m\)
and the same obstruction.
If such a positive foreign contact existed, its actual shadow
would force the same opposite to close an orthogonal component,
contrary to the original-farthest obstruction proved above.
Terminal four has foreign remainder \(<1/m\le1/5\), so no positive
spectator cost, and then negative time.
The other zero-own \(D/E\), internal-four and two-white roles
have their actual positive affine stars or cycles.
Their time is positive through the actual cap and every omitted
wall product is nonpositive.  Own contact is therefore positive.

Write \(q=\sqrt{D_i}\).  Here \(5/2<q\le3\).
The actual edge ratio needed for a loss \(\ell\) is
\[
 R_\ell(q)={1\over\ell}
                -{q-\ell\over q(2q-\ell)}.
\]
It increases with \(q\) and decreases with \(\ell\) on this
rectangle.  Thus
\[
 R_\ell(q)\le R_{.6}(3)=41/27<761/500.
\]
The finite-white residual formulas give the following ordinary
lower ratios, with side order at least five:
near-three/own at least five has
\((101/125+2/5)/\sqrt{3/5}\);
a \(D\)-position has
\((101/125+1/5)/\sqrt{2/5}\);
the \(E\), terminal-four and two-end-three rows are stronger.
All exceed \(761/500\).
The smaller owns have the same actual affine principals.
For the two exceptional roles WITH positive foreign cost,
the complete clearings give ratio
\(>227/(25\sqrt{35})>761/500\), or ratio square \(>121/35\).
These estimates retain the ENTIRE contacted spectator cost.
Hence no additional ordinary nonisolated row is weak.

Unlike the older \(3/2\) test, the present ratio threshold retains
near-own-four also at side order seven.  Up to reflection all
remaining such profiles are
\[
          v=\omega_i/\sqrt2+\omega_{i-1}/2,\qquad i=5,6,7 .
\]
The simple-three helper normalizes the cap to square two over the
actual field, with old product \(-1\); its own product is
\(-\sqrt2\), so \(\sqrt2\in K\).
At the embedding reversing \(\sqrt2\), the time numerators at
\(N=16\) are
\[
        126-88\sqrt2,\quad143-100\sqrt2,\quad154-108\sqrt2 .
\]
They are positive since \(\sqrt2<17/12\).
Their increments with \(N\) are
\((i-1)(3-2\sqrt2)>0\).
The conjugate new-root Schur complement is the negative of this
numerator divided by \(N\), contradicting positive definiteness.
This removes all three side orders and both orientations.

The remaining nonisolated possibility is pure:
\[
 c=Ap-\gamma\omega_i-\omega_e/2,\qquad
 e=1\text{ or }N-1,\quad Q=0,
\]
with own five, finite own at least seven, or dotted.
The exact helper-side identity is
\[
 {(\gamma+1/(2m))^2\over(m+1)/(2m)}-(4\gamma-2)
 ={(\gamma-1-1/(2m))^2\over(m+1)/(2m)}\ge0.
\]
The ratio threshold and native-five gap therefore give
\[
             202/125<t=2\gamma<3139/1458<54/25.
\tag{Astrip}\label{eq:arithAhelperphysical}
\]

For isolated-within-\(A\) contact, retain the whole foreign
projection \(v_S\) and write
\[
 c=Ap-\gamma\omega_i-v_S,\quad Q=\|v_S\|^2,\quad
 u=(1-Q)/\gamma^2,\quad A^2=\gamma^2(D_i-u).
\]
Here \(Q<1\) by finite-white positivity.
When \(Q>0,u>1\), the actual positive clearing has loss at most
\[
 {1\over\sqrt{D_i}+\sqrt{D_i-1}}+\sqrt{u-1}
 <.2+.383<\ell_N .
\]
This uses \(Q\ge1/4\), \(\gamma\ge\varphi/2\), and
\(u-1\le5-3\varphi<.383^2\).
At \(u=1\) the positive cleared vector is a forbidden feasible
future null.
The interval \(0<u<1\) is RETAINED, including \(Q>0\);
we do not charge a strong floor there.

For pure isolated owns, three, four, six and parallel have their
actual affine/null principals.
Own five has the actual \(H_4\) projected-point loss
\(<25/76+5/24<.6\).
For own seven, the increasing cubic
\(8a^3-4a^2-4a+1\) gives \(.9<a<.902\).
The actual positive \(A_2\) core gives
\[
 h_i>\sqrt{D_i-a^{-2}}-\sqrt{4/3-a^{-2}},
\]
and hence loss
\[
 <{100/81\over2.6+2.37}+\sqrt{17/150}
 <.25+.337=.587 .
\]
Pure own eight is arithmetically impossible:
\(4\gamma^2=2+\sqrt2\), and its opposite conjugate would require
\(D_i(2-\sqrt2)<4\), whereas
\(D_i\ge55/8>4+2\sqrt2\).
Every embedding of this minimal field extends to the actual
totally real field, so this contradiction applies there.
For all finite owns at least nine, the exact singleton finite
projection gives
\(h_i>\sqrt{D_i-1}-\tan(\pi/m)\).
Since \(\tan(\pi/9)<1386/3727<2/5\), the loss is \(<.6\).
This covers the infinite finite-own tail, not a finite label sample.
All positive finite and negative-norm projections just used are
positive sums of ACTUAL roots; their selected signs, positive
time and all omitted walls are retained as in the companion's
actual-point lemmas.

\subsection{Every pure weak helper is golden over the actual field}
\label{app:arithgoldenfield}

Normalize old \(A\) roots to square two.
If \(r\) is a \(K\)-root on the actual cap line and its simple-three
helper is \(a_e\) of square two, then
\(r^2=2(r,a_e)^2\).  Rescaling by \(-(r,a_e)^{-1}\in K\)
therefore makes the cap square two with helper product \(-1\).
Its own product is \(-t\).
The path cycle through old simple-three bonds makes \(t\) integral.
This compares only connected root lines; no disconnected norm
squareclasses are equated.

Every nonidentity embedding of \(\mathbb Q(t)\) extends to \(K\).
Strict conjugate positivity of the full old-component Schur
complement gives
\[
                    D_it_\sigma^2+2C_{ie}t_\sigma+D_e<4 .
\]
After reversing \(i\le N/2\), the cheaper costs for negative and
positive \(t_\sigma\) are respectively
\[
 Q_-(s)=2is^2+4s+2-2(is+1)^2/N,\qquad
 Q_+(s)=2is^2+2-2(is-1)^2/N .
\]
At the smallest \(N=16,i=5\), their values at \(-4/5,12/25\)
exceed four by \(3/40,59/1000\), respectively.
The costs at those endpoints increase with \(i\) and \(N\);
their values at zero are below four.
Convexity of the quadratics therefore gives the full strict strip
\[
                         -4/5<t_\sigma<12/25 .
\tag{Agoldenstrip}\label{eq:arithAgoldenstrip}
\]
This covers EVERY nonphysical conjugate, not chosen ones.

If the degree is at least five, the nonzero integral norm has
magnitude \(<(54/25)(4/5)^4<1\).
At degree four or three its magnitude is below two.
Any positive small conjugate would make the magnitude below one,
so all small conjugates are negative and their absolute product
is \(1/t\).
At degree four the integral trace lies in \((-1,0)\):
its lower bound is \(t-12/5>-1\), and its upper bound
\(t-3t^{-1/3}<0\) follows from \((54/25)^4<27\).
At degree three the trace lies in \((0,1)\), since
\(t-8/5>0\) and \(t-2/\sqrt t<1\);
the latter follows from
\((54/25)(29/25)^2<4\).
Both degrees are impossible.
At degree two the integral trace is one or two.
Trace two makes the nonzero norm \(t(2-t)\) lie in \((-1,1)\).
Trace one forces \(t<9/5\) and \(t(1-t)\in(-2,0)\);
the norm is \(-1\), and \(t=\varphi\).
Degree one leaves only \(t=2\), the actual parallel pair with
a feasible future null.
Thus every pure weak helper has \(\gamma=\varphi/2\).

The old component and this foreign-free cap have a physically
Lorentzian Gram over \(\mathbb Q(\varphi)\).
If \(K\) were a proper extension, a nonphysical real embedding
fixing that subfield would preserve this indefinite principal
Gram, impossible.
Consequently the helper forces the ACTUAL field
\(K=\mathbb Q(\varphi)\), rather than a convenient overfield.

\subsection{The complete golden packing calculation}
\label{app:arithgoldenpacking}

For a golden cap reversed to \(5\le i\le N/2\), its time
numerator \(M=2NA^2\) is
\[
 M_i^{\rm near}=i(N-i)-(N+1)+(i+2)(N-i)\varphi,\quad
 M_i^{\rm far}=i(N-i)-(N+1)+i(N-i+2)\varphi .
\]
These formulas retain BOTH helper ends.
At \(N=18\), conjugate positivity excludes the midpoint nine
and opposite-end indices six through nine.
The endpoint costs are \((173-84\varphi)/36\) and
\((197-99\varphi)/36\), both greater than one, and the
opposite-end cost is concave in \(i\).
At \(N=17\), only far seven is excluded, by
\((170-84\varphi)/34>1\) using \(\varphi<34/21\).
At \(N=16\) retain every orientation.
The resulting COMPLETE necessary tables are
\[
\begin{array}{c|c|r}
 N&\text{type and }M&\mathrm N(M)\\ \hline
16&\mathrm{near}\ 5:\ 38+77\varphi&-1559\\
 &\mathrm{near}\ 6:\ 43+80\varphi&-1111\\
 &\mathrm{near}\ 7:\ 46+81\varphi&-719\\
 &\mathrm{near}\ 8:\ 47+80\varphi&-431\\
 &\mathrm{far}\ 5:\ 38+65\varphi&-311\\
 &\mathrm{far}\ 6:\ 43+72\varphi&-239\\
 &\mathrm{far}\ 7:\ 46+77\varphi&-271\\ \hline
17&\mathrm{near}\ 5:\ 42+84\varphi&-1764\\
 &\mathrm{near}\ 6:\ 48+88\varphi&-1216\\
 &\mathrm{near}\ 7:\ 52+90\varphi&-716\\
 &\mathrm{near}\ 8:\ 54+90\varphi&-324\\
 &\mathrm{far}\ 5:\ 42+70\varphi&-196\\
 &\mathrm{far}\ 6:\ 48+78\varphi&-36\\
 &\mathrm{far}\ 8:\ 54+88\varphi&-76\\ \hline
18&\mathrm{near}\ 5:\ 46+91\varphi&-1979\\
 &\mathrm{near}\ 6:\ 53+96\varphi&-1319\\
 &\mathrm{near}\ 7:\ 58+99\varphi&-695\\
 &\mathrm{near}\ 8:\ 61+100\varphi&-179\\
 &\mathrm{far}\ 5:\ 46+75\varphi&-59 .
\end{array}
\]
At \(N=16\), near and far eight are reflections at the same
central index, so one row suffices.

Why do these finite tables restrict coexistence?
Subtracting a cap's \(K\)-valued projection onto the FULL old
forest leaves a nonzero vector on the SAME one-dimensional
time line over \(K\).
The caps have no foreign component.
Thus any two time coefficients have ratio in \(K\);
their \(M\)-ratio must be a square IN \(K\), and its rational
norm must be a rational square.
This argument retains all spectators.

At \(N=16\) and \(18\), the positive norm magnitudes in each
table are pairwise coprime nonsquares; their ratios are therefore
not rational squares.  These are direct small integer checks of
the displayed complete tables.
At \(N=17\), different norm squareclasses \(-1,-19,-179\)
cannot mix.
In class \(-1\), write the four entries as
\(42\varphi^3,18\varphi^5,14\varphi^5,6\varphi^7\).
After squares are removed the rational coefficient classes are
\(42,2,14,6\).
Every pair ratio has an odd valuation at three or seven;
none is a rational square or five times a rational square.
Those are exactly the positive rational squares in
\(\mathbb Q(\sqrt5)\), as expansion of
\((a+b\sqrt5)^2\) shows.
For the repeated \(-19\) class,
\[
 {M_{\rm far8}\over M_{\rm near6}}
   ={25+33\varphi\over76}={83+33\sqrt5\over152}
\]
has norm \(1/16\).
A square would require
\(a^2=121/304\) or \(45/304\), both with odd denominator
valuation at nineteen.
Thus different types cannot coexist, including the repeated
norm classes where a norm-only test would be insufficient.

At a fixed index the actual opposite is unique.
Reflected copies of the same noncentral type would have
\[
 -2c_L\cdot c_R=
 { (N-2i)(i+(i+2)\varphi)-(N+2)\over N}
\]
for near helpers, and the same formula with \(i-2\) replacing
\(i+2\) for far helpers.
Their simple-three helper bonds and old path make this
parameter integral by the invariant cyclic criterion, even for
a dotted mutual.
But the constant numerators for \(i=5,6,7\) at \(N=16\) are
\(12,6,-4\); for \(i=5,6,7,8\) at \(17\) they are
\(16,11,2,-11\); and at \(18\) they are \(20,16,8,-4\).
The far rows have the same respective constants.
None is divisible by \(N\), contrary to
\(\mathcal O_{\mathbb Q(\varphi)}=\mathbb Z[\varphi]\).
The central \(N=16\) row cannot occupy two different indices.
Hence at most ONE central index has a golden helper.
This is an all-native-five packing statement, independent of
which of those caps are actually weak at the new loss.

\subsection{Sum every retained isolated foot in one reference}
\label{app:arithsummedfeet}

The intervals \(0<u<1\) retained earlier cannot be declared
strong by the older \(.61\) estimate.
For such a cap clear the entire spectator projection by its
positive actual-root combination.  The actual vector
\[
 \eta=c+v_S+\gamma u n_i
      =\gamma\{\sqrt{D_i-u}\,p-(\omega_i-un_i)\}
\]
has square \(-\gamma^2u(1-u)\), positive time, zero contacted
spectator products, negative or zero \(c,n_i\) products, and
all other actual products nonpositive.
Its normalization is a genuine point of the whole chamber.
Original maximality gives
\[
 x_A\cdot(\omega_i-un_i)
 \ge\sqrt{D_i-u}-\sqrt{u(1-u)}.
\]
If \(\lambda_i=\sqrt{D_i}-\sqrt{D_i-u}+\sqrt{u(1-u)}\)
denotes its loss, then
\[
 \lambda_i\le a_i u+\sqrt{u(1-u)},\qquad
 a_i={1\over\sqrt{D_i}+\sqrt{D_i-1}}.
\]
The rational constants needed below are
\[
\begin{array}{c|ccc}
N&a_i<a&\text{positive reference coefficient }\beta_i\ge b&
 \ell\\ \hline
16&.1982&.022&.601\\
17&.1955&.0205&.6\\
18&.194&.018&.6 .
\end{array}
\]
The \(a\) bounds are positive-square comparisons at
\(D_{\min}=55/8,120/17,65/9\).
For every \(0<u<1\),
\[
 a u+\sqrt{u(1-u)}-\frac b2(2u-u^2)<\ell.
\tag{Aabsorb}\label{eq:arithAabsorb}
\]
For \(u\le1/2\) use \(a/2+1/2\), giving
\(.5991,.59775,.597\).
For \(u\ge1/2\), use \(2u-u^2\ge3/4\) and
\(\max(au+\sqrt{u(1-u)})=(a+\sqrt{1+a^2})/2\).
Radical upper bounds \(1.0195,1.019,1.019\) give respectively
\(.6006,.5995625,.59975\), all strictly below their \(\ell\).

Let \(H\) be a concave reference with available nonnegative
coefficients \(\beta_i\), positive ones at least \(b\):
\[
 w_j=(H_j-H_{j-1})/\sqrt2,\quad
 \beta_i=w\cdot n_i,\quad
 w=\sum_i\beta_i\omega_i,\quad
 E_H=\|w\|^2=\sum_i\beta_iH_i .
\]
Require \(H_i+\ell<\sqrt{D_i}\) at its positive central
coefficients.  Omit the possible golden index.
If \(S\) is the set of retained isolated indices with
\(\beta_i>0\), define ONE vector
\[
                       Z=w-\sum_{i\in S}\beta_i u_i n_i .
\]
Sum all actual point pressures BEFORE estimating its norm:
\[
 x_A\cdot Z\ge E_H-\sum_{i\in S}\beta_i(\lambda_i-\ell).
\]
Because \(w\cdot n_i=\beta_i\) and
\(n_i\cdot n_j\le0\) for \(i\ne j\),
\[
 \|Z\|^2\le E_H-\sum_{i\in S}\beta_i^2(2u_i-u_i^2).
\]
It follows that
\[
 e_A\ge2x_A\cdot Z-\|Z\|^2
 \ge E_H+\sum_{i\in S}\beta_i
       \{\beta_i(2u_i-u_i^2)-2(\lambda_i-\ell)\}>E_H.
\tag{Asummed}\label{eq:arithAsummed}
\]
Strictness follows from~\eqref{eq:arithAabsorb} if \(S\ne\varnothing\),
and from the strict available constraints otherwise.
Adjacent isolated indices and positive foreign costs are allowed.
No nonadjacency is assumed for this new isolated set, and no
independently optimized energies are added.

The references used by the endpoint arguments are
\[
\begin{array}{c|c|c|c}
N&H_5,H_6,\ldots\text{ to midpoint}&E_H&
 \text{maximum internal golden omission loss}\\ \hline
16&2.021,2.137,2.205,2.227&.8354522&.000576\\
17&2.056,2.186,2.269,2.310&.8708972&.00055225\\
18&2.087,2.228,2.324,2.381,2.399&.9037838&.00050625 .
\end{array}
\]
Reflect \(H_{N-i}=H_i\), use linear end tails, and at \(N=17\)
take \(H_9=H_8\).
The positive curvatures have exactly the lower bounds \(b\) above.
All floor inequalities are positive-square comparisons:
the new heights are the old \(.61\) heights shifted by
\(.61-\ell\), so \(H_i+\ell\) is unchanged.
A single internal omission sets its coefficient zero by neighbour
averaging and increases the adjacent coefficients.
An end omission replaces the first tail by the chord \(0\)--\(6\),
or its reflection.  These are changes in ONE reference.
In particular the full and omitted references preserve the
summed-foot argument whenever their positive central coefficients
have the displayed lower bounds; subsequent endpoint grafts must
check that condition again rather than assume it.

\subsection{A golden end hole pays for its own endpoint}
\label{app:arithgoldenend}

For later proper-component consumers we record how the exceptional
end omission is closed.  Suppose its central index is five,
reversing if necessary, and put \(\gamma=\varphi/2<.81\).
For either helper end \(e\), with its adjacent old node \(j\),
the actual finite \(H_4\) correction is
\[
 b_c=\gamma n_5+\tfrac23n_e+\tfrac13n_j,\qquad
 \eta=c+b_c,\quad
 \eta^2=\tfrac23-\gamma^2<.111^2 .
\]
The actual positive finite projected point has pure spatial
displacement \(v-b_c\).
Its normalization and original maximality give
\(x\cdot(v-b_c)\ge A-.111\).
Since \(x\cdot b_c>0\), the helper-end signed coordinate satisfies
\[
                  x_{\rm end}>
             \sqrt2(A-.111-\gamma h_5)>.49 .
\]
Both orientations are covered: the smaller opposite-end time
squares, at \(N=16,17,18\), are
\[
 (38+65\varphi)/32,\quad(42+70\varphi)/34,\quad
 (46+75\varphi)/36.
\]
They give \(A>2.115,2.136,2.155\).
Weakness gives \(h_5<2.022,2.057,2.088\).
Using \(\sqrt2>1.4\), the corresponding signed-coordinate lower
bounds are \(.512652,.502362,.493808\).
The nearer-end time is larger.
Here the signed coordinate is \(x_1\) at helper one and
\(-x_N\) at helper \(N-1\).
No strong floor at the missing index five has been inferred.

Use the chord \(0\)--\(6\) at the omitted end.
At the helper end the resulting reference is flat on \(L=6\)
coordinates in the near case and \(L=5\) in the far case,
with signed mean \(.25<b_0<.3\).
Replace that block by
\[
             .49,\quad
             {Lb_0-.49\over L-1},\ldots,
             {Lb_0-.49\over L-1}.
\]
Its sum is unchanged.  The latter coordinates lie in
\((.19,.262)\); the next old coordinate is below \(.101\).
Thus concavity persists, the new central boundary coefficient
is \(>.062>b\), and all other positive central coefficients
keep their previous lower bound.
The new endpoint coefficient is positive and is paid by the
actual endpoint pressure just proved.
The one-reference energy increases by
\[
                {L\over L-1}(.49-b_0)^2>.04332 .
\]
The pre-graft chord energy is exactly
\[
 E_H-\frac12\{H_5^2/5+(H_6-H_5)^2-H_6^2/6\}.
\]
At \(N=17,18\) it is respectively
\(12569249/15000000\) and \(13079276/15000000\);
at \(N=16\) it exceeds \(.8008\).
These direct rational formulas give, after the graft and
simultaneous absorption~\eqref{eq:arithAsummed}, the conservative
component bounds
\[
                 e_{A_{15}}>.8348,\qquad
                 e_{A_{16}}>.8703,\qquad
                 e_{A_{17}}>.9032 .
\]
If the single golden hole is internal instead, subtract the
internal losses in the reference table:
the energies are \(>.8348762,.87034495,.90327755\).
If there is no hole they are stronger.
The three conclusions therefore include every omission location,
all retained isolated feet, and arbitrary finite spectators.
They are the shared component estimates; no historical
ambient-dimension consumer has been imported.



